\begingroup
% Tighter columns + slightly taller rows for readability
\setlength{\tabcolsep}{3.8pt}
\renewcommand{\arraystretch}{1.08}

\begin{table*}[ht]
  \centering
  \small
  \caption{%
  The $R^2$ evaluation metrics for the fitted scaling laws, holding out separate dimensions of generalization: the largest model sizes $N$, most training tokens $D$, most compute $C$, and unique multilingual training mixtures $M$. In both monolingual and multilingual settings, we average $R^2$ across languages=[EN, FR, RU, ZH, HI, SW], including [ES, DE] for the multilingual setting. \textbf{We find \ourscl{} outperforms prior work in Monolingual and Multilingual settings.} We ablate the use of the terms for $D_{\mathrm{other}}$ and transfer languages ($\sum_{\mathcal{K}_{t}}D_i$).}
  \begin{tabular}{@{}>{\centering\arraybackslash}m{3mm}|@{\hspace{2pt}}l@{\hspace{6pt}}c@{\hspace{4pt}}c@{\hspace{4pt}}c@{\hspace{4pt}}c@{\hspace{4pt}}c@{}}
    \toprule
    & \textsc{Scaling Laws} & \textbf{$R^2$} & \textbf{$R^2(N)$} & \textbf{$R^2(D)$} & \textbf{$R^2(C)$} & \textbf{$R^2(M)$} \\
    \addlinespace[1pt]
    \cmidrule(l{4pt}r{4pt}){2-7}

    \rowgroup{3}{Mono.}
    & \basicscl{} \citep{hoffmann2022training}                  & 0.94 & 0.68 & 0.94 & 0.90 & -- \\
    & \textsc{Data-Constrained Scaling Law} \citep{muennighoff2024scaling}       & 0.93 & 0.78 & 0.93 & 0.88 & -- \\
    & \textbf{\underline{[Ours]}} \oursclabb{} \textsc{($D_t$ only)}              & 0.92 & \textbf{0.88} & 0.91 & 0.88 & -- \\
    \addlinespace[2pt]
    \cmidrule(l{4pt}r{4pt}){2-7}

    \rowgroup{5}{Multi.}
    & \basicscl{} \citep{hoffmann2022training}                             & 0.64 & -0.99 & 0.72 & 0.66 & 0.61 \\
    % \cmidrule(l{4pt}r{4pt}){2-7}
    & \textsc{Multilingual Scaling Law} \citep{he2024scaling}& 0.67 & -0.65 & 0.73 & 0.67 & 0.70 \\
    & \textbf{\underline{[Ours]}} \oursclabb{} \textsc{($D_t$ only)}                       & 0.70 & -0.75 & 0.80 & 0.72 & 0.64 \\
    & \textbf{\underline{[Ours]}} \oursclabb{} \textsc{($D_t+D_{\mathrm{other}}$)}         & \textbf{0.98} & \textbf{0.89} & \textbf{0.97} & \textbf{0.97} & 0.66 \\
    & \textbf{\underline{[Ours]}} \oursclabb{} \textsc{($D_t+D_{\mathrm{other}}+\sum_{\mathcal{K}_t}D_i$)} & \textbf{0.98} & \textbf{0.89} & \textbf{0.96} & \textbf{0.98} & \textbf{0.82} \\
    \bottomrule
  \end{tabular}
  \label{tab:evaluations}
  \vspace{-2mm}
\end{table*}
\endgroup


% \vspace{-1mm}
\section{Adaptive Scaling Laws for Monolingual \& Multilingual Settings}
\label{sec:multilingual-scaling}
% \vspace{-1mm}




\noindent\textbf{Challenges with existing scaling laws for multilingual modeling.} In this section, we discuss our attempts to precisely model different monolingual and multilingual pretraining experiments.
Scaling laws for English are typically based on Chinchilla \citep{hoffmann2022training}---which we refer to as the \basicscl{} (\basicsclabb).
\citet{hoffmann2022training}'s two power-law terms (\Cref{eq:scl-base}), allows us to separate how loss scales model size $N$ and data size $D$.
$E$ is the irreducible loss (representing the entropy of natural text), ($A$, $B$) are learned coefficients and and ($\alpha$, $\beta$) are scaling exponents for for model size and data size respectively.
Beyond English, languages may have different scaling laws \citep{arnett2025language}, and language resources are often limited, requiring a \dcsl{} (\dcslabb) \citep{muennighoff2024scaling} to account for diminishing returns after multiple epochs.
However, \dcslabb{} requires ample data both before and after one epoch to accommodate its two-stage fitting process.
For high-resource languages (English, French, Chinese), it is costly to collect ample data beyond one epoch, and for low/mid-resource languages (even Hindi, Hebrew, and Swahili) it can be very difficult to collect sufficient observations below one epoch.\footnote{In MADLAD-400 these languages have $<0.3\%$ English tokens. At standard batch sizes, even frequently evaluating (every $1k$) isn't enough to collect many (sometimes any) observations before 1 epoch is reached.}
As such, even for monolingual modeling, we require a scaling law that is data repetition-aware, and robust to different collection allowances for $(N,D)$.

For modeling multilingual mixtures, monolingual scaling laws can be used (either with $D$ representing the total or target language data), or \citet{he2024scaling} introduce \mscl{} (\msclabb), which expresses the loss for target language $t$ using $N,D$ and the sampling ratio for the target languages' language family in the training mixture.
However, each of these solutions only accounts for one of: the sampled target language data $D_{t}$, the total data altogether $D_{t} + D_{\mathrm{other}}$, or a set of the target and likely positive transfer languages $D_t+\sum_{i} D_i$.
We posit that a better model would separate and learn to weight these independent contributions.
In summary, none of the existing multilingual solutions account for (a) multi-epoch data repetition, or (b) the cross-lingual transfer effects beyond the target's language family.

% \clearpage

\begin{equation}
\mathcal{L}(N,\mathcal{D}_{\mathrm{eff}})
  = E \;+\; \frac{A}{N^{\alpha}} \;+\; \frac{B}{\mathcal{D}_{\mathrm{eff}}^{\beta}}
  \label{eq:scl-base}
\end{equation}
% \\[4pt]
\begin{equation}
\mathcal{D}_{\mathrm{eff}}
  = \underbrace{\textcolor{d_target}{\mathcal{S}_{\lambda_t}\!\left(D_t;U_t\right)}}_{\text{Monolingual}}
   \;+\; \underbrace{\textcolor{d_transfer}{\sum_{i\in\mathcal{K}} \tau_i \,\mathcal{S}_{\lambda_i}\!\left(D_i;U_i\right)}}_{\text{Transfer Languages}}
   \;+\; \underbrace{\textcolor{d_other}{\tau_{\mathrm{other}}\, \mathcal{S}_{\lambda_{\mathrm{other}}}\!\left(D_{\mathrm{other}};U_{\mathrm{other}}\right)}}_{\text{Other Languages}}
  % \qquad \tau_j,\tau_{\mathrm{other}}\ge 0
  \label{eq:scl-deff}
\end{equation}

\begin{equation}
\mathcal{S}_{\lambda}(D;U) = 
\begin{cases}
D, & D \le U \,\,\,\,\,\, (\le1\,\text{epoch}) \\
U\left[1 + \frac{1-\exp(-\lambda(D/U-1))}{\lambda}\right], & D > U \,\,\,\,\,\,(>1\,\text{epoch})
\end{cases}
\label{eq:scl-sat-piecewise}
\end{equation}

\noindent\textbf{The \ourscl{}.}
To resolve these challenges, we introduce the \ourscl{} (\oursclabb), a simpler variant of \dcslabb{} that is repetition-aware, introduces fewer additional parameters (for the monolingual variant), is fit in one stage, and is adaptable to an open-ended number of languages that would benefit from an explicit cross-lingual term.
In \Cref{eq:scl-base} the core scaling law formula includes the standard parameters governing the irreducible loss and $N,D$ scaling: $E,A,B,\alpha,\beta$.
Its effective data exposure term $D_{\mathrm{eff}}$ (\Cref{eq:scl-deff}) unpacks data sources into three terms: a $\textcolor{d_target}{\text{Monolingual}}$ term for the target language data $D_t$, an optional $\textcolor{d_transfer}{\text{Transfer Language}}$ term for up to $|\mathcal{K}_{t}|$ languages we wish to learn independent transfer coefficients for $\tau_{i}$, and an $\textcolor{d_other}{\text{Other Languages}}$ remainder term that sums all training tokens not already accounted for in the first two terms $D_{\mathrm{other}}=D_{\mathrm{tot}}-D_t-\sum_{\mathcal{K}_{t}}D_i$.
The latter two terms are optionally added for multilingual modeling, and $\mathcal{K}_{t}$ can be adapted as any combination of the languages with presumed highest positive transfer to the target language, or languages with the greatest representation in the experiment training mixture.
We found the latter was most effective and selected the $|\mathcal{K}_{t}|=3$ most highly co-sampled languages along target language $t$ across mixtures.
For each term, we apply a saturation function (\Cref{eq:scl-sat-piecewise}), ensuring a smooth decay on the effective data for each subsequent epoch where the training tokens have surpassed that languages' unique tokens $U$.
The new parameters introduced over standard scaling laws are: the repetition parameter $\lambda$, which is shared across each data source, the transfer weights $\tau_i$ for each language in $\mathcal{K}_{t}$ (which are initialized from language transfer scores derived in \Cref{sec:lang-transfer}), and the transfer weight $\tau_{\mathrm{other}}$ for the remainder of language tokens.


\noindent\textbf{Research Question:} \emph{How do scaling laws differ by language, and by monolingual vs multilingual training mixtures?}

% \sayna{ok - this first question seems to be the most important question in the paper (hence teaser figure is also about this) - however based on my limited prior knowledge about this topic I think the following text answers the question at surface and doesn't get into the "why" - so it talks about compute tax (the curves shift up) but maybe we should decompose this tax by looking at the parameters? does the tax come from a worse scaling exponent (alpha/beta) or a higher E? i think the latter for the multilingual setting rather than a degradation of the scaling exponent - example for that is the irreducible loss for Chinese which increases from 1.08 in the mono setting to 1.25 in the regular setting which suggests that while a multilingual vocabulary introduces a higher error due to capacity competition, the actual rate of improvement with more data and parameters remains robust - although the unimax setting uses positive transfer to lower both the exponent and the loss for most languages - this is turning tax into a net benefit }
\noindent\textbf{Monolingual scaling behavior is consistent in form, and multilingual variants exhibit a variable-sized compute-efficiency tax.}
\Cref{fig:monolingual-scaling} compares optimal scaling trajectories across six-variable resourced languages, under three regimes: \emph{monolingual vocabulary + monolingual training}, \emph{multilingual vocabulary + monolingual training}, and \emph{multilingual vocabulary + unimax training}.
Across languages, the monolingual curves are nearly parallel, indicating similar exponents and comparable returns to additional data and parameters (see full law parameters in \Cref{tab:mono-scaling-laws}).
Both multilingual vocabulary and \unimax{} training shift the frontier upwards, evidencing a \emph{compute-efficiency tax} relative to the monolingual-vocabulary/monolingual-training setting.
This is most pronounced for English, indicating it benefits less from language transfer than other languages do from English.
For Hindi and Swahili, the right tail bends upward, consistent with diminishing returns from severe data repetition after many epochs.
These qualitative trends motivate a single functional form that remains stable across languages while accounting for vocabulary/training-regime effects.

\noindent\textbf{Research Question:} \emph{How well can our scaling laws capture unique monolingual constraints, and complex multilingual cross-lingual transfer dynamics?}

\noindent\textbf{\oursclabb{} outperforms prior scaling laws, with a more robust fit across held-out axes in monolingual and multilingual settings.}
\Cref{tab:evaluations} evaluates fitted laws by holding out the largest model sizes \(N\), token ranges \(D\), compute \(C=6ND\), and held-out language mixtures $M$, reporting \(R^2\) averaged over available languages.
In the monolingual setting, all scaling laws perform comparably across most dimensions of generalization.
However, when generalizing to the largest size of model, laws that model data repetition outperform those that don't.
\oursclabb{} provides the strongest generalization to greater model sizes: \(R^2(N)=\mathbf{0.88}\) versus \(0.68\) (\basicsclabb{}) and \(0.78\) (\dcslabb{}), respectively.

In the multilingual setting, monolingual scaling laws can only observe their target data tokens, and as such perform adequately, but not well.
In particular, they are unable to generalize to the largest models $R^2(N)$ at all.
This is also the case for \msclabb{}, although it obtains better generalization to unseen language mixture in $R^2(M)=0.69>0.64$ by virtue of modeling the whole language family.
By separating the sources of data, \oursclabb{} is able to significantly outperform the other laws across all dimensions, frequently achieving $>0.9$ fit.
However, our ablation of the data terms shows that only by adding the $\textcolor{d_transfer}{\text{Transfer Language}}$ term is the model able to outperform \msclabb{} and achieve a better multilingual mixture generalization $R^2(M)=0.82$.
In summary, \oursclabb{} offers a simple framework to adaptively model cross-lingual transfer, and enables reliable extrapolation across out-of-sample dimensions in both monolingual and multilingual settings---addressing a critical gap where existing approaches fall short.